% Clase 15 — circuito RC con interruptor de dos posiciones y la curva de carga.
% Posición A: la batería carga el condensador. Posición B: la batería queda
% fuera y el condensador se descarga por R. Q es la carga de la placa a la que
% llega la corriente I, de modo que dQ/dt = I.
% Derecha: Q(t) = B C (1 - e^{-t/RC}), en unidades de B C y de RC.
\begin{tikzpicture}[scale=1]

  \begin{scope}[line width=0.9pt, color=figblue]
    % bateria y riel inferior
    \draw (0,0) to[battery1] (0,2.4);
    \draw (0,2.4) -- (1.2,2.4);
    \draw (0,0) -- (5.2,0);
    % posicion B: puentea la bateria
    \draw (1.2,1.5) -- (1.2,0);
    \fill (1.2,0) circle (2pt);
    % bornes del interruptor
    \draw[fill=white] (1.2,2.4) circle (2pt);
    \draw[fill=white] (1.2,1.5) circle (2pt);
    % polo y palanca (en A)
    \draw (1.26,2.36) -- (2.05,1.95);
    \fill (2.05,1.95) circle (2pt);
    \draw (2.05,1.95) -- (2.5,1.95);
    \draw (2.5,1.95) to[R] (4.5,1.95);
    \draw (4.5,1.95) -- (5.2,1.95);
    \draw (5.2,1.95) to[C] (5.2,0);
  \end{scope}
  \node[figlbl, anchor=east] at (-0.45,1.2) {$\mathcal{B}$};
  \node[figlblaux] at (1.2,2.75) {A};
  \node[figlblaux, anchor=east] at (1.05,1.5) {B};
  \node[figlbl, figblue] at (3.5,2.5) {$R$};
  \node[figlbl, figblue, anchor=east] at (4.8,0.97) {$C$};
  \node[figlblaux, text=figred, anchor=west] at (5.55,1.3) {$+Q$};
  \node[figlblaux, text=figblue, anchor=west] at (5.55,0.65) {$-Q$};
  \draw[figvecres] (4.35,2.3) -- (5.05,2.3);
  \node[figlbl, text=figred] at (4.7,2.62) {$I$};

  % curva de carga
  \begin{scope}[xshift=8.4cm, yshift=-0.2cm]
    \begin{axis}[
      emfig, width=6.2cm, height=4.6cm,
      xmin=0, xmax=5.2, ymin=0, ymax=1.2,
      xtick={0,1,2,3,4,5}, ytick={0,0.5,1},
      xlabel={$t/RC$}, ylabel={$Q/(\mathcal{B}C)$},
    ]
      \addplot[figaux, domain=0:5.2, samples=2] {1};
      \addplot[figblue, domain=0:5.2, samples=80] {1-exp(-x)};
    \end{axis}
  \end{scope}

\end{tikzpicture}
